\section{Background and Motivation: Learned Cache Eviction}
\label{sec:background}

At each full-cache miss, a caching system must evict one resident object before admitting the requested one. Classical policies such as LRU operationalize this choice through simple recency structure, while learned or learning-augmented approaches attempt to predict which eviction will lead to lower future miss cost. In practice, these approaches differ in feature sets, prediction targets, and evaluation methodology.

The main difficulty for a benchmark is that victim selection is structured, not i.i.d. A model does not act on one object in isolation; it acts within a decision set whose size, feature spread, and ambiguity vary across traces and capacities. A useful benchmark should therefore preserve decision boundaries instead of flattening them away.

This observation motivates the core \lafc\ design choice: represent supervision at the candidate-within-decision level. Each row corresponds to one candidate victim at one full-cache eviction decision, and the label records the finite-horizon counterfactual miss cost of forcing that eviction under a fixed continuation policy. This representation allows both pointwise learning and within-decision evaluation, while keeping the supervised object aligned with the actual cache-eviction choice.

\begin{figure}[t]
  \centering
  \fbox{\parbox{0.9\columnwidth}{\centering Placeholder for benchmark task overview figure:\\candidate rows $\rightarrow$ decision view $\rightarrow$ pairwise sample}}
  \caption{Planned overview of the three benchmark task views derived from candidate-level supervision.}
  \label{fig:task-overview}
  \Description{Placeholder box for a benchmark-overview figure showing candidate rows flowing into a decision view and a pairwise sample.}
\end{figure}

\todo{Add a compact illustrative example of one cache decision with 3--5 candidate victims and one optimal set.}
